% Pursuing Stacks (cotes 134-2 to 134-8): archive cote 134-3, site page 84 (PDF page 85), typescript page 84.
% typescript. Draft compared: ps2.tex:4373-4445.
% AI-transcribed by Claude Opus 5.5 (claude-opus-5-5, adaptive thinking, effort high), Message Batches API, batch ps1 (msgbatch_01Rv5hCaNTtUt4JHx1yqoHPL), 2026-11-12; image review recorded in records/CHECK_DECISIONS_134_3_4.tsv.
\scanpage{134-3/84}
\mspage{84}

(iii) \xout{A\hand{\ensuremath{^{\wedge}}} admit} A satisfies the two conditions \inserted{[besides $A\neq\emptyset$]}:\xout{T 2) and T 3)}
T 2) The product \inserted[typed]{in A\hand{\ensuremath{^{\wedge}}}} of any two elements in A is aspheric.
T 3) A\hand{\ensuremath{^{\wedge}}} admits a homotopy interval (which is equivalent to
saying that A is a local test category).

\inserted{Of course (ii)$\Longrightarrow$(iii).} The condition (iii) can be expressed in terms of the topos \hand{$\mathcal{A}$} = A\hand{\ensuremath{^{\wedge}}}
by saying that this is a totally aspheric and locally modelizing topos,
which implies already (as we saw by the \ensuremath{\overset{\hand{\vee}}{\mathrm{C}}}ech computation of cohomology)
that the topos is aspheric, and therefore modelizing,\inserted{ i.e. A a test category,} which is just (ii)\hand{.}
\inserted{Thus (ii)$\Longleftrightarrow$(iii).} Total asphericity of \hand{$\mathcal{A}$} , and the \xout{fact}\inserted[typed]{property} that \hand{$\mathcal{A}$} be a local modelizer, are
expressed respectively by the two neat conditions T 2) and T 3), neither
of which implies the other even for an A\hand{\ensuremath{^{\wedge}}}, with A a category with final
object\deleted{\unsure{s}}. \deleted{\unsure{In}}

For expressing in more explicit terms condition (i), and check equi-%
valence with the two equivalent conditions (ii),(iii), we need to \xout{assume}
admit that the canonical functor from (Cat) to it\over[typed]{'}{s} localization (Hot)
commutes with binary products\hand{$^{(36)}$.} \over{t}{T}his being so\hand{,} the condition of commutation
of A\hand{\ensuremath{^{\wedge}}}\hand{\ensuremath{\to}}(Hot) with binary products can be expressed by the \inserted[typed]{more concrete} condition that
for F,G in A\hand{\ensuremath{^{\wedge}}}\xout{\unsure{n}}, the canonical map in (Cat)
\[ (8) \qquad i_A(FxG)\hand{\ensuremath{\to}} i_A(F)xi_A(G) \qquad \text{i.e.} \qquad A_{/FxG}\hand{\ensuremath{\to}} A_{/F}xA_{/G} \]
is a weak equivalence. In fact, the apparently stronger condition\xout{s}
that the maps (8) are aspheric will follow , because the usual criterion
of asphericity for a map in (Cat) shows that it is enough \inserted{for this} to check that
(8) is a weak equivalence when F and G are representable by elements
a and b in A, which also means that $A_{/axb}$ is aspheric namely axb in A\hand{\ensuremath{^{\wedge}}}
aspheric - which is nothing but condition T 2) in (iii), \xout{This shows}
i.e. total asphericity of A\hand{\ensuremath{^{\wedge}}}. On the other hand, the condition that A
be a weak test category means that the elements $j_A(C)$ in A\hand{\ensuremath{^{\wedge}}}\hand{,} for C in
(Cat) admitting a final object, are aspheric, or (what amounts to the
same when A\hand{\ensuremath{^{\wedge}}}\xout{x} is totally aspheric) that \xout{\unsure{that is in}} \inserted[typed]{they are} aspheric over the final
object $e_{A^{\hand{\wedge}}}$ i.e. that the products $axj_A(C)$ are aspheric, which also \xout{mean}
means A is a local test category. Thus (i) \xout{impl} is equivalent to (ii).

The equivalent conditions (i) to (iii) are expressed by saying that
A is a \uline{strict test category} , or that \xout{the}\hand{$\mathcal{A}=$} A\hand{\ensuremath{^{\wedge}}} is a \uline{strictly modelizing}
\uline{topos} , or also that it is an\xout{strict} \uline{elementary} \uline{strict} \uline{modelizer} (when
emphasizing the topos \hand{$\mathcal{A}$} should be of the type A\hand{\ensuremath{^{\wedge}}}\hand{/} indeed). This presumably
will turn out to be the more important among the three notions\hand{,} of a \xout{weak}
``test category'' and the weak and strict variant. The two less stringent
notions however seem interesting in their own right\over{.}{,} \over{T}{t}he notion of a weak
test category mainly (at present) because it turns out that the category
